\section{Overview}
\label{sec:ch2-overview}

% Write your methods overview here.



\section{Exploratory Data Analysis}
\label{sec:exploratory-data-analysis}





\section{Model Specification}
\label{sec:model-specification}

In frequentist survival analysis, a frailty model is a type of survival model that accounts for unobserved heterogeneity among individuals or groups. This involves introducing a random effect, known as the frailty term, which captures the variability in risk that is not explained by observed covariates. The frailty term can be thought of as a latent variable that represents the individual's risk of death. In the Bayesian context this is formulated as a hierarchical model, where the frailty term is treated as a random variable with its own prior distribution. The model can be specified as follows:

Let \(i\) represent the individual patient, \(j\) represents the given ICU unit (\texttt{first\_careunit}), and \(t_{i,j}\) is the time-to-event (censoring) for patient \(i\) in ICU unit \(j\). The hazard function for patient \(i\) in ICU unit \(j\) can be expressed as:

\begin{equation}
    h_{i,j}(t) = h_0(t) \exp(X_{i,j} \beta + \alpha_i)
\end{equation}

Where \(h_0(t)\) is the baseline hazard function, representing the death over time for "average" patient, \(X_{i,j}, \beta\) are the fixed effects representing patient-level covariates (e.g., age, comorbidities) while \(\beta\) represents the log-hazard ratios for these predictors, and \(\alpha_i\) is the log-hazard cluster indicator for ICU unit \(j\). 

\subsection{Priors}


